\newglossaryentry{glossar}{
    name        = Glossar,
    description = {Laut Wikipedia ist ein Glossar \enquote{eine Liste von Wörtern mit beigefügten Bedeutungserklärungen oder Übersetzungen.}}
}

\newacronym[
    shortplural = {Abk.en},
    longplural  = {Abkürzungen},
    plural      = {Abk.en}
]{abk}{Abk.}{Abkürzung}

\newacronym[
    shortplural = {Power-on-Reset},
    longplural  = {Power-on-Reset-Schaltung},
    plural      = {PoR}
]{por}{PoR}{Power-on-Reset-Schaltung}

\newacronym[
    shortplural = {ICs}, 
    longplural  = {integrierte Schaltkreisen}, 
    plural      = {ICs}
]{ic}{IC}{integrierte Schaltkreise}

\newacronym[
    plural = {MOSFET-Transistoren}
    shortplural = {MOSFETs}
    longplural  = {Metall-Oxid-Halbleiter-Feldeffekttransistoren}, 
]{mosfet}{MOSFET}{Metall-Oxid-Halbleiter-Feldeffekttransistor}

\newacronym[
    shortplural = {CMOS}
    longplural  = {komplementären Metalloxid-Halbleitertransistoren},
]{cmos}{CMOS}{komplementärer Metall-Oxid-Halbleitertransistor}

\newacronym[
    longplural  = {n-Kanal Metalloxid-Halbleitertransistoren},
]{nmos}{NMOS}{n-Kanal Metall-Oxid-Halbleitertransistor}

\newacronym[
    longplural  = {p-Kanal Metalloxid-Halbleitertransistoren},
]{pmos}{PMOS}{p-Kanal Metall-Oxid-Halbleitertransistor}

\newacronym[
]{tcr}{TCR}{spezifische Temperaturkoeffizient}

\newacronym[
    shortplural = {OTAs},
    longplural  = {Operational-Transconductance-Amplifiers}
]{ota}{OTA}{Operational-Transconductance-Amplifier}

\newacronym[
    shortplural = {PVTs},
]{pvt}{PVT}{Prozess, Spannung und Temperaturbedingungen}

\newacronym[
]{mc}{MC}{Monte-Carlo}